% Integrable regular parts may be dropped in all local coordinate comparisons.
% Requires amsmath, amssymb, amsthm.  Labels are prefixed invg:.
% Suggested bibliography keys: IncomingExact, IncomingTransport, Latt2011.
\section{Which logarithmic singularities have no coordinate defect?}
\label{invg:sec}

The nonlinear coordinate formulas in the two incoming reports give a
finite algorithm for each specified singular germ.  The fifth research
question in the Exact Identities report asks for an intrinsic
classification of the germs whose defect vanishes under every coordinate
tangent to the identity.  We answer that question, and more generally
allow any prescribed order of tangency.  The necessary condition is as
important as the sufficient one: logarithmic degrees from different pole
orders cannot conceal a defect for every coordinate.

The change-of-variables mechanism for logarithmic finite parts is
classical; see, for example, L\"att's Theorem~2 in
\cite{Latt2011}.  The theorem below is an exact invariant-subspace
classification and a specialization to the Stieltjes hierarchy relative
to the audited incoming material.  No general historical-priority claim
is made.

\subsection{Conventions and the inherited residue formula}

Work locally on the positive side of the origin.  A logarithmic
meromorphic germ has the unique singular expansion
\begin{equation}\label{invg:germ}
 h(x)=\sum_{q=1}^{R}\sum_{n=0}^{L}c_{q,n}x^{-q}\log^n x+h_{\rm reg}(x),
\end{equation}
where the omitted part is locally integrable and is a finite sum of
analytic germs times nonnegative powers of $\log x$.  Coefficients may be
real or complex.  Define $\operatorname{FP}_x h$ by taking the constant
term of its cutoff integral against a smooth test function.  Positive
powers of the cutoff logarithm and negative powers of the cutoff are
removed.  Use scalar distribution pullback, so that ordinary functions
pull back by composition.

For an orientation-preserving analytic coordinate $f$, put
\[
 H_f(t)=f(t)/t,\qquad \lambda_f(t)=\log H_f(t),\qquad
 C_f[h]=f^*(\operatorname{FP}_x h)-\operatorname{FP}_t(h\circ f).
\]
The residue law proved in the incoming reports is
\begin{equation}\label{invg:residue}
 \langle C_f[h],\phi\rangle
 =\sum_{q,n}\frac{c_{q,n}}{n+1}
    \operatorname*{Res}_{t=0}
     f(t)^{-q}\lambda_f(t)^{n+1}\phi(t)\,dt.
\end{equation}
It is enough to keep the finite Taylor jet of the smooth test function
inside the residue.  Integrable terms create no discrepancy.  Equivalently,
for a single monomial,
\begin{equation}\label{invg:monomial}
 C_f[x^{-q}\log^n x]
 =n!\sum_{j=0}^{q-1}\frac{(-1)^j}{j!}
 [z^{n+1}t^{q-1-j}]H_f(t)^{z-q}\,\delta_0^{(j)}.
\end{equation}
This is the same coefficient law, because
$[z^{n+1}]H_f^{z-q}=H_f^{-q}\lambda_f^{n+1}/(n+1)!$.

For completeness, the analytic justification is short.  Continue
$x_+^{z-q}$ after subtracting the test-function Taylor polynomial through
degree $q-1$.  Only the term of degree $q-1$ has denominator $z$ near
zero.  After pullback, its numerator depends on $z$ through $H_f^{z-q}$.
The constant term of the cutoff coefficient omits the quantity
$n![z^{n+1}]$ of that numerator, whereas the meromorphic coefficient
retains it.  Every other Taylor denominator is nonzero at $z=0$; the
remainder is normally integrable there.  This proves
\eqref{invg:monomial}, including its sign and factorials.

\subsection{The complete classification}

For an integer $d\ge1$, write
\[
 \mathscr G_d=\{f:f(0)=0,\quad f(t)=t+O(t^{d+1})\}.
\]
These are germs with identity Taylor jet through degree $d$.

\begin{theorem}[Classification by pole order and logarithmic degree]
\label{invg:classification}
For the germ \eqref{invg:germ} and an integer $d\ge1$, the following
conditions are equivalent.
\begin{enumerate}
\item $C_f[h]=0$ for every $f\in\mathscr G_d$.
\item $C_{f_a}[h]=0$ for every sufficiently small real $a$, where
      $f_a(t)=t(1+a t^d)$.
\item $c_{q,n}=0$ whenever $q>d(n+1)$.
\end{enumerate}
Equivalently, if $h=\sum_{n=0}^{L}a_n(x)\log^n x$, the pole order of
$a_n$ is at most $d(n+1)$ for every $n$.  For $d=1$, the exact condition
is $\operatorname{poleord}(a_n)\le n+1$.
\end{theorem}

\begin{proof}
The first condition implies the second.  To prove that the third implies
the first, note that $\lambda_f(t)=O(t^d)$ for $f\in\mathscr G_d$.
The $q,n$ summand of the residue integrand in \eqref{invg:residue} has
order at least $d(n+1)-q\ge0$.  It is analytic at zero, so its residue
against any smooth test function is zero.

It remains to prove necessity; this is not a consequence of estimating
the order of one monomial alone.  For $f_a=t(1+a t^d)$,
\eqref{invg:monomial} becomes
\begin{equation}\label{invg:quadratic-family}
 C_{f_a}[x^{-q}\log^n x]
 =n!\sum_{k=1}^{\lfloor(q-1)/d\rfloor}
 \frac{(-1)^{q-1-dk}a^k}{(q-1-dk)!}
 [z^{n+1}]\binom{z-q}{k}\,\delta_0^{(q-1-dk)}.
\end{equation}
The term at $k=0$ vanishes because $n+1\ge1$.
For a fixed power $a^k$ and derivative $\delta_0^{(j)}$, the pole order
is uniquely determined by $q=dk+j+1$.  Therefore different pole orders
cannot cancel this coefficient.

Fix $q$ and let $K=\lfloor(q-1)/d\rfloor$.  The coefficient of
$a^k\delta_0^{(q-1-dk)}$ vanishes for $k=1,\ldots,K$ only if
\begin{equation}\label{invg:triangular}
 \sum_{n=0}^{k-1}c_{q,n}n![z^{n+1}]\binom{z-q}{k}=0,
 \qquad k=1,\ldots,K.
\end{equation}
This is triangular in $c_{q,0},\ldots,c_{q,K-1}$, and its diagonal
coefficient at $c_{q,k-1}$ is
\[
 (k-1)![z^k]\binom{z-q}{k}=\frac1k\ne0.
\]
Induction on $k$ gives $c_{q,0}=\cdots=c_{q,K-1}=0$.
Those indices are exactly the pairs satisfying $q>d(n+1)$.
This proves the third condition and completes the equivalence.
\end{proof}

A useful consequence is that the least required tangency order can be
read directly from the singular expansion:
\begin{equation}\label{invg:least-tangency}
 d_{\min}(h)=\max_{c_{q,n}\ne0}
              \left\lceil\frac q{n+1}\right\rceil.
\end{equation}
For a germ with no singular term, every coordinate is harmless.  For a
nonzero singular germ, \eqref{invg:least-tangency} is the least positive
integer $d$ such that matching the identity jet through degree $d$
removes every coordinate defect.

\begin{corollary}[Exact dimension of the obstruction space]
\label{invg:rank}
Within pole order at most $R$ and logarithmic degree at most $L$, the
codimension of the invariant subspace is
\begin{equation}\label{invg:rank-formula}
 \sum_{q=1}^{R}\min\!\left(L+1,\left\lfloor\frac{q-1}{d}\right\rfloor\right).
\end{equation}
The one-parameter family in \eqref{invg:quadratic-family}, with its
polynomial dependence on $a$ retained, detects all these obstructions.
\end{corollary}
\begin{proof}
Theorem~\ref{invg:classification} sets precisely the listed independent
coordinates $c_{q,n}$ to zero.  Equation~\eqref{invg:triangular} proves
that their obstruction functionals are independent.
\end{proof}

The reference to the whole parameter family matters: a single numerical
choice of $a$ can combine coefficients belonging to different pole orders.
The theorem uses an exact polynomial identity in $a$.

\subsection{Sharp Stieltjes thresholds for any order of tangency}

Use the usual convention
$\gamma_n(x)=x^{-1}\log^n x+\gamma_n(1+x)$.  Let
$T_{n,r}=\operatorname{FP}_x\gamma_n^{(r)}(x)$ be the finite part of
the pointwise argument derivative.  This object is distinguished from
$D^r(\operatorname{FP}_x\gamma_n)$, whose harmonic contact term was
computed in the incoming reports.

\begin{theorem}[Exact threshold for Stieltjes argument derivatives]
\label{invg:stieltjes-threshold}
For $r,n\ge0$ and $d\ge1$,
\[
 f^*T_{n,r}=\operatorname{FP}_t\bigl(\gamma_n^{(r)}\circ f\bigr)
 \quad\text{for every }f\in\mathscr G_d
\]
holds if and only if
\begin{equation}\label{invg:stieltjes-criterion}
 d\ge r+1\qquad\text{or}\qquad
 n\ge r+\left\lfloor\frac rd\right\rfloor.
\end{equation}
When $1\le d\le r$, set $K=\lfloor r/d\rfloor$, $e=r-dK$, and
$n_*=r+K-1$.  The first excluded index has the exact defect
\begin{equation}\label{invg:stieltjes-top}
 C_{f_a}[\gamma_{n_*}^{(r)}]
 =\frac{(-1)^e n_*!}{e!K!}a^K\delta_0^{(e)},
 \qquad f_a(t)=t(1+a t^d).
\end{equation}
Thus every threshold in \eqref{invg:stieltjes-criterion} is attained.
\end{theorem}

\begin{proof}
The singular part is
\begin{equation}\label{invg:stieltjes-singular}
 \gamma_n^{(r)}(x)
 =(-1)^rr!x^{-r-1}
 n![z^n]\left\{\prod_{j=1}^{r}\left(1-\frac zj\right)e^{z\log x}\right\}
 +\gamma_n^{(r)}(1+x).
\end{equation}
Every elementary symmetric coefficient of
$1,1/2,\ldots,1/r$ is nonzero.  Hence its smallest logarithmic degree is
$\max(n-r,0)$, and the coefficient at that degree is nonzero.
The classification theorem gives the exact condition
\[
 r+1\le d\bigl(\max(n-r,0)+1\bigr).
\]
If $d\ge r+1$ this holds for all $n$.  Otherwise it requires $n\ge r$
and is equivalent to the second inequality in
\eqref{invg:stieltjes-criterion}.

At $n=n_*$, the smallest logarithmic degree is $K-1$, with coefficient
$n_*!/(K-1)!$ in \eqref{invg:stieltjes-singular}.  Every higher
logarithmic degree is invariant.  In
\eqref{invg:quadratic-family}, only $k=K$ survives for this smallest
degree, and its derivative index is $e$.  Its coefficient is
$(-1)^ea^K/(e!K)$.  Multiplication by $n_*!/(K-1)!$ proves
\eqref{invg:stieltjes-top}.
\end{proof}

For $d=1$ this recovers the incoming sharp threshold $n\ge2r$ and its
coefficient $(2r-1)!a^r\delta_0/r!$.  It additionally gives, for example,
\[
 C_{t+a t^3}[\gamma_3^{(3)}]=-6a\delta_0',
 \qquad
 C_{t+a t^3}[\gamma_5^{(4)}]=60a^2\delta_0,
\]
while $\gamma_4^{(3)}$ and $\gamma_6^{(4)}$ have no defect under any
coordinate $f(t)=t+O(t^3)$.

\subsection{Consequences for coincident Stieltjes products}

The invariant space is linear but is generally not closed under
pointwise multiplication: $x^{-1}$ is invariant under every unit-tangent
coordinate, whereas $x^{-2}$ is not.  The classification nevertheless
provides a useful exact product criterion.

\begin{corollary}[Products with at most one digamma factor]
\label{invg:stieltjes-products}
Let $m_1,\ldots,m_s$ be nonnegative integers, at most one of which is
zero.  Then the specified pointwise-product finite part satisfies
\begin{equation}\label{invg:stieltjes-products-formula}
 f^*\operatorname{FP}_x\!\left(\prod_{j=1}^{s}\gamma_{m_j}(x)\right)
 =\operatorname{FP}_t\!\left(\prod_{j=1}^{s}\gamma_{m_j}(f(t))\right)
 \qquad (f\in\mathscr G_1).
\end{equation}
In particular, all powers $\operatorname{FP}_x\gamma_m(x)^s$ with
$m\ge1$ have this invariance.
\end{corollary}
\begin{proof}
Choose the singular term from a nonempty subset $S$ of the factors.
Before expanding the remaining analytic factors, the resulting term is
\[
 x^{-|S|}\log^{\sum_{j\in S}m_j}x
             \prod_{j\notin S}\gamma_{m_j}(1+x).
\]
Its logarithmic degree is at least $|S|-1$.  Taylor expansion of the
analytic factor can only decrease the pole order.  Thus every resulting
singular monomial satisfies $q\le n+1$, which is exactly the criterion
of Theorem~\ref{invg:classification} for $d=1$.
\end{proof}

The restriction on the number of zero indices cannot be replaced by a
large total Stieltjes index.  If $f(t)=t+a t^2+O(t^3)$ and
$\gamma_m=\gamma_m(1)$, direct use of the same classification gives
\begin{align}
 C_f[\gamma_0^2]&=a\delta_0,\label{invg:product-counter1}\\
 C_f[\gamma_0^2\gamma_1]&=
             \left(a\gamma_1+\frac{a^2}{2}\right)\delta_0,
             \label{invg:product-counter2}\\
 C_f[\gamma_0^2\gamma_m]&=a\gamma_m\delta_0\qquad(m\ge2).
             \label{invg:product-counter3}
\end{align}
Indeed, the mixed singular--singular--regular term in the last two
products is $\gamma_m x^{-2}$.  The fully singular term is
$x^{-3}\log^m x$: it contributes $a^2\delta_0/2$ when $m=1$ and
nothing when $m\ge2$.  All other terms meet the invariant criterion.
These are identities for a specified finite part of the ordinary
product; they do not postulate a product of distributions.

\paragraph{Verification.}
The accompanying exact script compares the residue construction with
\eqref{invg:quadratic-family}, computes the rank of the finite
obstruction map, and checks the Stieltjes criterion and its sharp
coefficient through several pole and tangency orders.  These are
independent algebraic audits of signs and finite formulas.  The
all-order conclusions follow from the proofs above.
